% SPDX-FileCopyrightText: 2026 Bart van de Lint
% SPDX-License-Identifier: AGPL-3.0-or-later
\section{How the model works}
\label{sec:details}

This appendix gives the formulas, calibrations and sources behind each step of \cref{sec:model}, in the same order.

\subsection{Hull, propeller and engines}
\label{sec:det_ships}

\paragraph{Calm water} The resistance $R_\mathrm{cw}(V)$ follows \citet{holtrop1982} and \citet{holtrop1984}\tbd{ (Holtrop 1984 read in part)}: the friction of a flat plate of the same wetted area by the ITTC-57 line times a form factor $1+k_1$, the friction of the appendages, wave-making, the resistance of the bulb and the transom, and a correlation allowance to full scale. The waterplane coefficient is $C_{WP}=0.763\,(C_P+0.34)$ from the prismatic coefficient $C_P$, a fit whose source we could not trace; the calibration below absorbs its error on the standard classes.

One factor $f_R$ per standard class (\cref{tab:classes}) multiplies $R_\mathrm{cw}$ so that the calm-water power delivered to the propeller at the service speed equals $\lambda_s P_\mathrm{MCR}\eta_\mathrm{tr}/(1+s)$: the source's rated power $P_\mathrm{MCR}$ at the load $\lambda_s$ at which it makes the service speed, through the shaft, without the sea margin $s$ (\cref{app:calm_calibration}). The coaster's source power is a fleet average rather than a design point, so the coaster and the scaled family keep $f_R=1$, and the family's engine is rated from its own calm-water power with any kite stowed, since it must make the service speed with the kite down.

\paragraph{Wind on the ship} In the apparent wind of speed $V_a$ and angle $\beta_a$ off the bow at \SI{10}{m}, the force along the ship and the side force are $X = \tfrac12\rho_a V_a^2 A_\mathrm{L} C_X$ and $Y = \tfrac12\rho_a V_a^2 A_\mathrm{L} C_Y$, with $\rho_a$ the air density and the loading functions of \citet{blendermann1993}
\begin{equation}
C_X = \frac{-C_{D,l}\cos\beta_a}{1-\tfrac{\delta}{2}\left(1-\tfrac{C_{D,l}}{C_{D,t}}\right)\sin^2 2\beta_a}, \qquad
C_Y = \frac{C_{D,t}\sin\beta_a}{1-\tfrac{\delta}{2}\left(1-\tfrac{C_{D,l}}{C_{D,t}}\right)\sin^2 2\beta_a},
\end{equation}
where $A_\mathrm{L}$ and $A_\mathrm{F}$ are the lateral and frontal areas above the waterline, $C_{D,t}$ is the drag coefficient in beam wind, $C_{D,l}=C_{D,l\mathrm{AF}}A_\mathrm{F}/A_\mathrm{L}$ with $C_{D,l\mathrm{AF}}$ the drag coefficient on the frontal area in head or stern wind, and $\delta$ the cross-force parameter. The side force heels the ship with a lever $\kappa_\mathrm{w}$ times the height of the centroid of $A_\mathrm{L}$, and turns it about a point $0.18\,L(\beta_a-\pi/2)$ aft of that centroid. Coefficients, areas and centroid are his for the nearest type, the bulk carriers taking the loaded tanker's, scaled to the class. A container class with more than his four tiers on deck adds the further tiers, each \SI{2.59}{m} high, over 77\,\% of its length. The resistance in \cref{eq:thrust} is $R_\mathrm{air}=-X$.

\begin{figure*}[tbp]
\centering
\plot[\textwidth]{wind_load}
\caption{Coefficients of the wind load on the ship itself over the apparent wind angle, on its lateral area, for one class of four types with their deck cargo, from the loading functions of \citet{blendermann1993}: (left) along the ship, positive forward, and (right) to the side. The other container ships and bulk carriers follow the Megamax and the Aframax tanker.}
\label{fig:wind_load}
\end{figure*}

\paragraph{Waves} The added resistance in irregular waves is the spectral integral of the regular-wave response,
\begin{equation}
R_\mathrm{aw} = 2 \int_0^\infty \frac{R_\mathrm{aw}(\omega,\mu)}{\zeta_a^2} S(\omega)\,\mathrm d\omega,
\end{equation}
where $R_\mathrm{aw}(\omega,\mu)$ is the added resistance in regular waves of frequency $\omega$, heading $\mu$ and amplitude $\zeta_a$, and $S(\omega)$ a JONSWAP spectrum, the standard form for a developing sea, with a peak enhancement factor of 3.3 \citep{hasselmann1973}. The regular-wave response is the semi-empirical SNNM formula as recalibrated by \citet{mittendorf2022}. In following seas we apply their energy-exchange rule to the whole motion term and interpolate that term between beam and following seas, where they interpolate only its amplitude coefficient. Because $R_\mathrm{aw}$ scales with $H_s^2$, the model tabulates $R_\mathrm{aw}/H_s^2$ over speed, peak period and heading.

\paragraph{Propeller and engines} A single propeller is an actuator disk with a viscous loss. Its diameter is half the draft on the scaled family, as on the Damen and Ferus Smit coasters, 0.6 of the design draught on the bulk carriers and tankers \citep{man2019mr}, and 0.74 on the container ships and the car carrier \citep{man2011}. For a net thrust $T$ the propeller delivers $T/(1-t)$ in an inflow $V(1-w)$, with the wake fraction $w=0.5C_B-0.05$ and the thrust deduction $t=0.6w$ of \citet{man2011}. Its efficiency is
\begin{equation}
\eta_0 = \eta_v \frac{2}{1+\sqrt{1+C_T}}, \qquad C_T = \frac{T/(1-t)}{\tfrac12\rho_w A_p V^2(1-w)^2},
\end{equation}
with the disk area $A_p$, the water density $\rho_w$ and a viscous efficiency $\eta_v=0.8$, so the delivered power is $P_D = T V (1-w) / [(1-t)\eta_0]$. The engine output is $P = P_D/\eta_\mathrm{tr} + P_\mathrm{rig} + P_\mathrm{hotel}$, with a transmission efficiency $\eta_\mathrm{tr}$ of 0.96 for the shaft and 0.90 for the electric chain, and the fuel power is $P/\eta$, with
\begin{equation}
\eta = \eta_\mathrm{peak}\left[1-k\,(1-\lambda)^2\right],
\end{equation}
where $\lambda$ is the load of the running units as a share of their rating, $\eta_\mathrm{peak}=0.44$ and $k=0.2$. A dual-fuel ammonia engine runs on the diesel cycle, and we take the diesel's efficiency for it, as no source we read gives another. On the scaled family the installed power is the calm-water output at \SI{10.5}{kn} plus a 15\,\% sea margin, at 90\,\% of the maximum rating.

\paragraph{Ammonia drive} The pilot share of 5\,\% is the low end of the 5 to 15\,\% expected \citep{mmmczcs2021fuel}. Liquid ammonia holds \SI{18.6}{MJ/kg} \citep{irena2021ship}. Its tank's volume \citep{mmmczcs2021fuel} amounts to \SI{600}{kg} of ammonia per cubic metre of hold, and it costs EUR~0.29 per kWh stored, USD~1.62 per kilogram \citep{kanchiralla2022}. No source we read gives the tank's mass, so we take the filled tank as twice its fuel, with bounds of 1.4 and 3.

\subsection{Hull and rudder forces}
\label{sec:det_balance}

With the symbols of \cref{sec:balance}, the hull's waterline length $L$, beam $B$, draught $T$ and block coefficient $C_B$, the lateral force slope of the hull, the side force per unit leeway angle, is, after Kijima et al. as given by \citet{kim2021},
\begin{equation}
Y'_\beta = \pi\frac{T}{L} + 1.4\,C_B\frac{B}{L},
\end{equation}
made non-dimensional with $\tfrac12\rho_w V^2 L T$. The side force acts at $x_\mathrm{h} = N'_\beta L/Y'_\beta = 2T/Y'_\beta$ forward of midship, from the yaw derivative $N'_\beta=2T/L$ of the same regression, fitted to 13 merchant-ship models on even keel with block coefficients of 0.52 to 0.84. The rudder at the stern, sized at the minimum area DNV recommends for a rudder directly behind a propeller, $A_r = TL/100\,[1+50C_B^2(B/L)^2]$ \citep{dnv_rudder}, cancels the yaw moment:
\begin{equation}
F_r = \frac{F_y x_\mathrm{h} - M_z}{-L/2 - x_\mathrm{h}}.
\end{equation}
The hull then carries $F_y+F_r$ at a leeway
\begin{equation}
\beta = \frac{F_y+F_r}{\tfrac12\rho_w V^2 L T Y'_\beta}
\end{equation}
and pays an induced drag $R_\beta$ of $(F_y+F_r)\beta$ for the cross-flow share $1.4C_BB/L$ of the slope and half of that for its potential-flow share $\pi T/L$, as for a slender wing. The rudder pays $R_r=F_r^2/(\tfrac12\rho_w V^2 A_r \pi \Lambda_r)$ with an effective aspect ratio $\Lambda_r$ of 1.6 \citep{kramer2022}. The heel is $\phi = (M_x + F_y T/2)/K_\phi$, with the righting stiffness $K_\phi = \Delta g\,GM$ from the displacement $\Delta$ and the metacentric height $GM$, the lever per radian of heel by which the hull rights itself, taken as 7.5\,\% of the beam.

\paragraph{Solving the trim} The engine output a setting needs is $P_D(T_\mathrm{req})/\eta_\mathrm{tr}+P_\mathrm{rig}$. It is evaluated at a few evenly spaced values of the trim, and Brent's method refines the best of them. A kite flown over wings or rotors has two trims, which the Nelder--Mead method refines together from the best scanned pair; one at a time would stall where the rigs together pull the ship and the propeller idles. This finds the optimum to within $5\times10^{-6}$ of the bare ship's engine output, where lists of settings miss it by up to 2.8\,\% for rotors, 1.9\,\% for the kite and 0.5\,\% for wings and show as wiggles in the rig forces over wind angle (\cref{fig:rig_polar}).

Each trim is solved against the calm-water resistance and the waves' added resistance together, and beyond the performance table's largest added resistance the propeller makes up the rest. Against calm water alone, a rotor with thrust to spare leaves the waves to the propeller, which took 0.8 points off the rotors' saving in cost per tonne-kilometre on the router's check cases (\cref{sec:routing}). Against lists of settings, the continuous solve raises the saving in cost per tonne-kilometre by 1.5 points on average for rotors, 0.7 for wings and 0.1 for the kite.

\subsection{Wing sails and rotors}
\label{sec:det_deck}

The largest units built set the split, \SI{40}{m} tall for wings \citep{alfalaval2026oceanbird} and \SI{35}{m} for rotors \citep{maritimeexecutive2024sohar}. The 75\,\% of the ship's length is our estimate of the deck a row can use, with a range of 65 to 80\,\%, from published layouts \citep{han2025,mboumba2022}; it caps the number of units only on the largest ships and elsewhere moves the thrust by less than 1\,\%. A maker offers 2 to 5 rotors on the bulk carriers and tankers \citep{anemoi2024brochure}, and the model lets the best rotor rigs reach \tbd{6 to 8} units on the Handysize, Ultramax and Kamsarmax, \tbd{13} on the Capesize, and \tbd{7, 10 and 16} on the MR, Aframax and VLCC tankers.

Wings stay at least 1.5 chords apart, so that wings turning about a quarter-chord pivot clear each other, also the closest spacing at which two-dimensional CFD of three wings in line found no further gain \citep{wang2024}. Rotors stay three diameters apart, the closest spacing measured in the wind tunnel \citep{bordogna2020}. The thrust of $n$ units $d$ chords or diameters apart is a factor
\begin{equation}
f_a = 1 - \left(1 - \frac{1}{n}\right) k_a\, e^{-d/d_0}
\end{equation}
of that of $n$ isolated units, averaged over wind angles. For rotors $k_a=0.38$ and $d_0=2.72$ diameters, fitted to wind-tunnel tests of two rotors \citep{bordogna2020} and CFD of three and four rotors \citep{han2025}. For wings $k_a=0.30$ and $d_0=2.86$ chords: CFD of sails held at one angle of attack gives $k_a=0.44$ \citep{malmek2024,xu2025}, of which we take back a third for sheeting each sail to the wind it meets \citep{vanreen2026}, and $d_0$ is our estimate. Both rigs take the apparent wind at their mid height from a power-law profile $U(z)=U_{10}(z/\SI{10}{m})^{0.05}$ fitted to the mean ERA5 model-level profiles over the study region.

A wing has an aspect ratio of 2 and a lift coefficient $C_L$ between 0 and 2.0, its trim, or it is stalled at an angle of attack $\alpha_s$ from \SIrange{25}{90}{\degree}, a second trim used running before the wind. Attached, its drag is
\begin{equation}
C_D = C_{D,0} + \left(k_s + \frac{1}{\pi e \Lambda_\mathrm{eff}}\right) C_L^2,
\end{equation}
with a profile drag $C_{D,0}$ of 0.02, a separation term $k_s=0.03$, a span efficiency $e=0.9$, and an effective aspect ratio $\Lambda_\mathrm{eff}$ 1.4 times the geometric one for the deck as end plate. Thrust and side force follow from the apparent wind angle $\beta_a$,
\begin{align}
F_x &= q A f_a\, (C_L\sin\beta_a - C_D\cos\beta_a),\\
F_y &= q A\, (C_L\cos\beta_a + C_D\sin\beta_a),
\end{align}
with the sail area $A$ and the dynamic pressure $q=\tfrac12\rho_a V_a^2$. Stalled, the force is normal to the wing with a coefficient $C_N=1.2$, so $C_L=C_N\sin\alpha_s\cos\alpha_s$ and $C_D=C_{D,0}+C_N\sin^2\alpha_s$.

A rotor has a height of six diameters. Its lift, drag and spin power coefficients over the spin ratio $\alpha=\Omega r/V_a$, with the angular speed $\Omega$ and radius $r$, are the ShipCLEAN polynomials, fitted to CFD and corrected with full-scale measurements \citep{tillig2020}. The spin power is $P_\mathrm{rig} = q A V_a C_{P,\mathrm{spin}}(\alpha)$, with the projected area $A$, height times diameter. The trim is the spin ratio, any value between 0 and 4, capped by a surface speed of \SI{47}{m/s}.

\paragraph{Rig size and deck space} The size search starts from three sizes spanning \SIrange{500}{16000}{m^2} for kites and a fifth to all of the deck capacity for wings and rotors, then adds two more, each at the vertex of the parabola in log area through the best size and its neighbours. Only on a container ship do wings and rotors cost cargo space: each unit blocks the on-deck stack of every \SI{2.5}{m} row its footprint covers, plus a metre of clearance, over one 40-foot bay and all tiers on deck. A rotor's footprint is its diameter. A wing sweeps a circle wider than its chord, so it must stand above the stacks or retract \citep{gcaptain2021vplp}; we raise it on a mast \SI{2.5}{m} wide, our estimate, which is its footprint. The kite's launch mast takes one metre.

\paragraph{Mass or volume limit} The cargo of a container ship or car carrier is $m_c=\min[(m_\mathrm{dwt}-m_\mathrm{rig}-m_\mathrm{stores}-m_\mathrm{tank})/(s_\ell\,m_u+m_\mathrm{tare}),\ N-N_\mathrm{rig}-V_\mathrm{tank}/V_u]\,s_\ell\,m_u$, with the deadweight $m_\mathrm{dwt}$, the masses of rig, stores and filled fuel tank, the tank volume $V_\mathrm{tank}$, $N$ the nominal capacity of \cref{tab:classes}, $N_\mathrm{rig}$ the slots the rig blocks, the cargo $m_u$ of a loaded unit, the tare $m_\mathrm{tare}$ and space $V_u$ of one unit, and the share $s_\ell$ of the units that are loaded, 0.812 for containers and 1 for car decks (\cref{app:cargo_units}). On the other ships a tank that does not fit in the engine room is assumed to fit in space the cargo leaves unused. No class changes group with the fuel store of any drive or route. The Megamax reaches its deadweight 5\,\% before its slots and becomes space-limited if a loaded TEU carries \SI{9.5}{t} of goods instead of \SI{10.2}{t}, and the Neo-Panamax becomes deadweight-limited at \SI{11.1}{t}, both within the range of \cref{tab:inputs}.

\paragraph{Fixed or retractable} Rotors on bulk carriers can be tilted, folded or moved on rails \citep{gcaptain2021seazhoushan,maritimeexecutive2024sohar,anemoi2026}, while one maker fits rotors on tankers fixed as standard \citep{anemoi_deploy}. A fixed rig prevents cargo handling on a general cargo or container ship, a geared bulk carrier loses \SI{2}{h} per call and a gearless one \SI{4}{h}, and tankers and car carriers lose no time. A ship in ballast floats 45\,\% of its draft higher, and where the top of a fixed rig then exceeds the lowest common clearance of the class, we take a share of its calls, 1 to 10\,\%, as lost for two days each.

\subsection{Kite}
\label{sec:det_kite}

\paragraph{Shape and aerodynamics} The design shape is a circular arc with a projected span 82\,\% of its flat span, a rounded planform and one cambered airfoil, NACA 4415 (\cref{fig:kite}). How the canopy deforms under load is neglected, and its losses enter through the flight drag below. The Vortex Step Method \citep{cayon2023}, not validated on a ram-air shape, uses section polars from NeuralFoil \citep{sharpe2025} with a drag increment of 0.015 for the cell openings and wrinkles. A projected aspect ratio of 3 is a flat aspect ratio of 3.76, and published values are lower, below 2 for the SkySails kites (\cref{app:kite_aspect}). \tbd{A projected aspect ratio of 2.5, the lowest with a usable polar on this shape, changes the cost saving on the three ships of the routing noise check (\cref{sec:checks}) by +0.9 to $-1.9$ points, with the flight drag held or recalibrated to the same glide ratio, about the size of that noise.} The SkySails sizes are read as flat and converted at 1.28 \citep{folkersma2020}.

\paragraph{Bridle and tether} The bridle has four chordwise rows of attachments, cascading to the control unit 0.7 spans below the canopy. Its lines carry the force limit at a working stress of \SI{0.4}{GPa}, and their drag adds a coefficient of $C_{D,\mathrm{tether}}\,\ell d/A$ for a total line length $\ell$. HMPE rope, a high-strength polyethylene fibre, of \SI{30.6}{mm} breaks at \SI{880}{kN}, so a design factor of 3, the usual start for a creep life of up to two years \citep{bosman2013,fechner2013}, gives this stress. We take the nominal diameter, not the worked-in diameter 15\,\% smaller of \citet{bosman2013}, which would make the rope denser than solid HMPE.

A quarter of the tether's drag acts at the kite \citep{vandervlugt2019}. Holding its length in flight, as on the SkySails ships \citep{erhard2013}, keeps it from bending over the drum under load. The smaller SkySails kite flew on 15.5 to 19.7 spans \citep[Fig.~18]{erhard2013}. \tbd{On the three ships of the routing noise check (\cref{sec:checks}), 15.5 spans change the cost saving by at most +0.3 points and 19.7 spans by +1.9 to $-1.0$.} Half the tether's weight hangs at the kite with the canopy, the bridle and a control unit of \SI{0.22}{kg} per kN of force limit, from the \SI{70}{kg} pod of a \SI{320}{kN} SkySails kite \citep{paulig2013}. That pod is powered through the tether \citep{maass2013}, and we leave this cable out.

\paragraph{Tether life} Creep governs a rope held under load at a fixed length \citep{bosman2013,humeau2018}. We sum its damage with Miner's rule, which adds up the fractions of life used at each load, over the hours each ship of \cref{sec:checks} flies at each tension on its routed trips. A power law gives about \num{11000} flight hours and a time--tension--temperature diagram about \num{24000}, and we take the geometric mean; the range of \cref{app:tether_life} moves the cost saving by up to 2.4 points.

Bending is charged per launch (\cref{eq:cycle_cost}). Pumped over a drum under load a tether fails after about \num{1100} flight hours \citep{bosman2013}, and bent over a fixed fairlead in every figure eight within tens of hours \citep{davies2015}, hence the fixed tether length and a tow point that turns with the tether. When the section at the tow point wears, reeling in a few centimetres moves fresh rope there, an estimate as no source gives the wear rate. On the three ships of \cref{sec:checks}, each with the life the creep laws give it, a thicker tether at \SI{0.3}{GPa} lowers the cost saving by \tbd{0.3 to 1.5} points and a thinner one at \SI{0.48}{GPa} changes it by \tbd{$-0.4$ to $+0.2$} points, within the routing noise, so we keep \SI{0.4}{GPa} as a stated choice.

\paragraph{Lift to drag and the upwind limit} Against the 8 to 8.5 of a rigid scale model in the wind tunnel \citep{poland2025}, kites reach about 4 in flight \citep{oehler2019}. Canopy deformation, the control unit and steering make up the difference, added as the flight drag coefficient of 0.08, which brings the kite with its bridle and tether to $G=3.9$ at the best angle of attack, from 8.4 for the canopy alone (\cref{fig:kite}). Line drag and length scale with the kite, so $G$ holds for every size.

The window edge lies at $\arctan G$ off downwind \citep{erhard2013}. Combined with the window edge at elevation $\theta$, $\cos\phi=\sin\varepsilon/\cos\theta$ with $\varepsilon=\arctan(1/G)$ \citep{bigi2018}, the force has a forward component on the ship only if the apparent wind at the kite comes from at least
\begin{equation}
\beta_\mathrm{min}=\arcsin\frac{\sin\varepsilon}{\cos\theta}
\label{eq:upwind_limit}
\end{equation}
off the bow, our own form of these two relations. The limit applies at the kite's height, so ship speed moves it in true wind angle. The force table respects it by construction, since each pattern is flown with this drag.

\paragraph{Canopy cloth} The canopy holds \SI{3.5}{m^2} of cloth per square metre of projected area in its two skins, ribs and diagonals, as paragliders do, plus \SI{0.05}{kg/m^2} of seams, tapes and reinforcements. The one full-scale kite we can check leaves \SI{0.19}{kg/m^2} beyond cloth and lines (\cref{app:canopy_mass}), and we keep \SI{0.05}{kg/m^2} as a stated choice. \tbd{At \SI{0.19}{kg/m^2} the cost saving on the three ships of the routing noise check (\cref{sec:checks}) changes by $-1.1$, $-0.6$ and +0.1 points.} A kite built for another force limit gets the cloth that keeps the same life under \cref{eq:uv_life}: over limits of 0.25 to \SI{2}{kN/m^2} the cloth weight grows with the limit to the power 0.27, and at the low bound of the wear exponent $b$ to the power 0.8, close to the upper bound \cref{sec:res_force} runs. Seams, tapes and reinforcements grow in proportion to the limit, while their cost stays in the making cost, so a kite built for twice the limit would pay about 5\,\% more for its canopy than we count.

\paragraph{Canopy price} The price per projected square metre (\cref{fig:canopy_price}) is
\begin{equation}
p_\mathrm{canopy} = f_\mathrm{m}\,(r_\mathrm{c}\,c_\mathrm{cloth} + c_\mathrm{make}),
\label{eq:canopy_price}
\end{equation}
the market factor $f_\mathrm{m}$ of 2 times the cost to make it: $r_\mathrm{c}$ square metres of cloth at a price $c_\mathrm{cloth}$ per sewn square metre, and the making cost $c_\mathrm{make}$ of sewing, bridle and hardware (\cref{app:canopy_price}). The goal of USD~150 keeps the sewing rate, the lines, the hardware and the factor. No mill has quoted the cloth, and sewing falls from 1.5 hours per flat square metre to 1.0, the top of what we infer from one year of US cargo parachute production, with no credit for a sewing robot. We keep the market factor, the factor of the other rigs, though the about 1.6 of airbags suggests it could fall (\cref{sec:costs}).

Sewing hours and cloth per square metre come from paragliders, whose cells, the air-filled chambers between the ribs, are several times narrower than a ship kite's. No source gives hours for either; if a ship kite needs fewer hours and less cloth, its cost saving rises by 0.2 to 1.0 points (\cref{app:seam_cloth}). Wright's law, under which the time per unit falls by a fixed ratio with each doubling of cumulative output, reaches the goal's 1.0 hour from a paraglider series' 1.5 hours in four of eight cases at \num{10000} canopies (\cref{app:learning_curve}), so the goal is a number for a mature series.

\paragraph{Quasi-steady flight} In the frame of the apparent wind, with $x$ downwind, $y$ crosswind and $z$ up (\cref{fig:kite_frames}), the kite's unit position on the sphere of radius $r$, the tether length, about the tow point is
\begin{equation}
\hat{\mathbf e}_r = (\cos\vartheta,\ \sin\vartheta\cos\phi,\ \sin\vartheta\sin\phi),
\end{equation}
with the angle $\vartheta$ off the downwind direction and the roll angle $\phi$ about it, zero horizontal and \SI{90}{\degree} straight up. The kite holds a heading $\psi$ in the tangent plane, measured from the direction $\hat{\mathbf e}_\vartheta$ towards the edge of the wind window, and flies through the air along $\hat{\mathbf h} = \cos\psi\,\hat{\mathbf e}_\vartheta + \sin\psi\,\hat{\mathbf e}_\phi$ at a speed $v_h$. With the wind velocity $\mathbf v_w$ at the kite, the apparent wind at the kite is $\mathbf v_a = (\mathbf v_w\cdot\hat{\mathbf e}_r)\,\hat{\mathbf e}_r - v_h\hat{\mathbf h}$ and its ground velocity $\mathbf v_k = \mathbf v_w - (\mathbf v_w\cdot\hat{\mathbf e}_r)\,\hat{\mathbf e}_r + v_h\hat{\mathbf h}$. For a kite of mass $m$, projected area $A$ and lift and drag coefficients $C_L$ and $C_D$, with the unit vector $\hat{\mathbf z}$ upwards, the force balance along and across the airspeed direction $\hat{\mathbf e}_a$ reads
\begin{align}
T_t\,\hat{\mathbf e}_r\cdot\hat{\mathbf e}_a &= \tfrac12\rho_a|\mathbf v_a|^2 A C_D - m(\mathbf a + g\hat{\mathbf z})\cdot\hat{\mathbf e}_a, \label{eq:kite_drag}\\
\left|\mathbf F_\perp\right| &= \tfrac12\rho_a|\mathbf v_a|^2 A C_L, \qquad \mathbf F = m(\mathbf a + g\hat{\mathbf z}) + T_t\hat{\mathbf e}_r, \label{eq:kite_lift}
\end{align}
where $\mathbf F_\perp$ is the part of $\mathbf F$ normal to $\hat{\mathbf e}_a$, $T_t$ the tether tension and $\mathbf a$ the acceleration of the turn, $v_h^2/R$ across the heading, and of the motion on the sphere, $|\mathbf v_k|^2/r$ towards its centre. \Cref{eq:kite_drag} gives the tension for a trial speed, and false position with the Illinois step solves \cref{eq:kite_lift} for $v_h$. Without mass and flying across the wind this reduces to Eq.~39 of \citet{loyd1980} with zero weight and the roll angle of his Eqs.~40--41, $T_t = \tfrac12\rho_a (\mathbf v_w\cdot\hat{\mathbf e}_r)^2 A \sqrt{C_L^2+C_D^2}\,(1+G^2)$.

A pattern is the heading law of the SkySails pattern generator \citep[Sec.~VIII]{erhard2013}. The kite holds the heading $+\psi_0$, turns through the edge direction to $-\psi_0$ when its roll angle passes $\phi_0+\phi_a$, and back when it passes $\phi_0-\phi_a$. Its yaw rate follows the turn-rate law $\dot\psi = g_s\,v_h\,\delta$ \citep[Eq.~7]{erhard2013}, so it turns on a fixed radius $R = 1/(g_s\delta)$ through the air, \SI{45}{m} for the kite of \SI{125}{m^2} projected, scaled with the span. Each leg flies straight at its heading, so a sweep runs along the roll angle: up and down where the kite flies at the side of the window, across it overhead. The model keeps these straight legs. Letting the heading also swing by up to \SI{45}{\degree} about $\pm\psi_0$ within each leg raises the best thrust found from the same 21 starts by at most \SI{2.6}{\percent} of the force limit, by \SI{1}{\percent} or more in at most three of 52 pairs of wind and course, and by at most \SI{0.5}{\percent} within \SI{30}{\degree} of the wind, for kites of \num{800} and \SI{1600}{m^2} and side-force weights from 0 to 3. The swinging patterns that gain are figure eights with reshaped turns, not sweeps along the edge. The tether force on the ship is the time average $\overline{\mathbf F}$ of $T_t\hat{\mathbf e}_r$ over the second settled cycle.

\paragraph{Pattern choice} For each apparent wind speed $V_a$ and ship course $\chi$ off the apparent downwind direction, a Nelder--Mead search, a direct search that needs no derivatives, picks the pattern that minimises
\begin{equation}
J = -\frac{F_\parallel - w\,|F_\perp|}{T_\mathrm{max}},
\label{eq:pattern_objective}
\end{equation}
where $F_\parallel$ and $F_\perp$ are the components of $\overline{\mathbf F}$ along and across the course, $T_\mathrm{max}$ the force limit and $w$ a weight on side force. The search varies the centre $\phi_0$ and half width $\phi_a$ of the roll angle, the heading $\psi_0$ and the angle of attack, up to \SI{2}{\degree} below stall, with the whole cycle above \SI{8}{\degree} of elevation. It starts from two figure eights and one sweep along the edge and from the best pattern of the next stronger wind. Where none of these flies it starts instead from the three best of a coarse grid of \num{3360} patterns over those four parameters, since a kite with little drag flies in only a small part of that space and the fixed starts can all miss it. Without the grid, zero flight drag lowers the cost saving of the Neo-Panamax from \tbd{9.0} to \tbd{5.3}\,\%; with it, zero flight drag raises that saving from \tbd{9.7} to \tbd{10.3}\,\%. The kite's air speed at each point of a pattern is found from a rising scan, since a heavy kite makes too little lift to hang at rest and flies once it moves; refusing such points would lower the cost saving of the best kite by \num{0.7} points on average. A different starting bracket for that solve moves one ship's cost saving by \num{-0.7} to \num{+1.0} points. The weight $w$ is the kite's trim, any value from 0 to 3; the search runs at seven weights and the forces are interpolated linearly in between. The apparent wind is taken at \SI{25}{\degree} of elevation on the tether, about \SI{270}{m} up for the base \SI{800}{m^2} kite (\cref{app:perf_grid}).

\paragraph{Kite held still} Besides flying a pattern, a kite can be held still: at a fixed point of the window with a fixed heading, its nose into the wind at its angle of attack, where lift, drag, weight and tether tension balance with no motion. This is \cref{eq:kite_drag,eq:kite_lift} with $\mathbf v_k = 0$ and $\mathbf a = 0$, so the air meets the kite at the wind $\mathbf v_w$ itself. Without mass it rests at $\vartheta = \arctan(C_L/C_D)$ off the downwind direction, and its weight draws it lower in the window. For a roll angle $\phi$ and an angle of attack, false position solves \cref{eq:kite_lift} for $\vartheta$. A Nelder--Mead search, started from the best of a grid of 19 roll angles and 7 angles of attack, picks the still kite with the least $J$ (\cref{eq:pattern_objective}), with the kite above the same \SI{8}{\degree}. Each cell of the table holds whichever of the best pattern and the best still kite has the lower $J$, and always the still kite where no pattern pulls the ship forward. A still kite at the side of the window pulls like a sail with the kite's glide ratio. So it is chosen in light wind from ahead of the beam, where it pulls more than the patterns the search finds, near the force limit in the strongest winds, and towards the wind, where a kite kept up costs only its drag. Dead upwind the still \SI{1600}{m^2} kite holds back the ship with \tbd{6, 14 and \SI{39}{kN}} at 6, 10.5 and \SI{18}{m/s}, against \tbd{85, 20 and \SI{43}{kN}} for the slowest sweep the search found.

The \SI{25}{m/s} flight limit is the cut-out wind of the SkySails Power land system \citep{skypower100_tib} and the top of the airborne range of the Kitepower Falcon \citep{kitepower2024}. Where no pattern pulls the ship forward the kite is held still, and the router keeps it up only while its drag costs less than a landing and a later launch. No other limit applies on the port approach. The launch limit is one Beaufort step below the Beaufort~8 the ship system was designed for \citep{noaa2009}, and the landing limit is the wind up to which the automated recovery of the \SI{100}{kW} land system is meant to work \citep{skypower100_tib}. A launch in a head wind may need the ship turned through \SI{180}{\degree} and back \citep{wiwo2008}; no source says how many do, so each launch carries a mean turn cost and the router does not turn.

\paragraph{Launch wind in a sea} A sea of significant wave height $H_s$, peak period $T_p$ and heading $\mu$ raises the lowest apparent wind for a launch by
\begin{equation}
\Delta V_\mathrm{sea} = \Delta V_\mathrm{ref}\,\frac{v_m(V, H_s, T_p, \mu)}{\bar v_{m,\mathrm{ref}}},
\label{eq:sea_wind}
\end{equation}
where $v_m$ is the significant vertical speed of the ship at its launch mast, twice the root-mean-square, from heave and pitch out of phase in long-crested waves with a JONSWAP spectrum, and $\Delta V_\mathrm{ref}=\SI{5}{m/s}$ the rise reported on the reference ship, whose $\bar v_{m,\mathrm{ref}}$ is the mean over headings at \SI{6}{m/s} in a sea of \SI{3.25}{m} and \SI{9}{s}. Heave and pitch come from the closed-form expressions for a box-shaped hull of \citet{jensen2004}, checked there against model tests of two fast ferries, a container ship and a VLCC, tabulated per metre of $H_s$ over speed, period and heading. Roll is left out, as a mast on the centreline barely rises with it. Since the scale is fixed on one ship, only the ratio of motions between ships matters, which a closed-form hull gives better than it gives their size.

\paragraph{Launch and landing} One launch and landing costs
\begin{equation}
c_c = (1+p_r)\left[c_\mathrm{crew} + p_t C_\mathrm{canopy} + p_l (C_\mathrm{canopy} + C_\mathrm{tether} + C_\mathrm{unit})\right] + c_\mathrm{turn} + \frac{n_b}{N_b} C_\mathrm{tether},
\label{eq:cycle_cost}
\end{equation}
with a chance of retry $p_r$ of 5\,\%, crew time $c_\mathrm{crew}$ of USD~85, a chance $p_t$ of a tear that ends the canopy, a chance $p_l$ of losing the kite in the water, and a mean turn cost $c_\mathrm{turn}$ of USD~130. Gradual wear from flogging and folding is in the canopy life. The last term is the tether's bending: $n_b=4$ bends (2 to 6) over drum and sheave, of which $N_b=\num{400000}$ end it, extrapolated from published laws for ropes of up to \SI{20}{mm} at higher tensions (\cref{app:tether_bending}). At the base the bending of a launch moves no cost saving by 0.01 points, and at the low bound with six bends it takes up to 0.5 points off. Launch and landing times are held at every kite size; reeling the whole tether at an assumed \SI{1.5}{m/s} would add about \SI{22}{min} each way at \SI{16000}{m^2}, which the upper ends of the two times (\cref{tab:inputs}) about reach.

\subsection{Lowest wind for flight and launch}
\label{sec:det_min_wind}

A kite in steady flight at elevation $\beta$, in an apparent wind $V_a$, with glide ratio $G=C_L/C_D$, sees the wind $V_a\cos\beta$ along the tether and its weight $mg\cos\beta$ in the tangent plane. Sweeping level across the wind, it rolls by $\psi$ so that $L\sin\psi=mg\cos\beta$. Its air speed is $V_a\cos\beta\sqrt{1+G^2\cos^2\psi}$, and the roll that makes most of it has $\sin^2\psi=(1+G^2)/(3G^2)$, so that
\begin{equation}
V_\mathrm{cross} = \sqrt{\frac{3\sqrt{3}\,mg}{\rho_a A\cos\beta\;\max_\alpha\left[C_L(1+G^2)^{3/2}/G\right]}},
\label{eq:min_cross}
\end{equation}
which for a large $G$ maximises $C_L^3/C_D^2$, the square of the endurance parameter of an aircraft \citep{echeverri2020part1}, beside relatives for a drag-power kite \citep{bauer2018} and a rolled kite \citep{luchsinger2013,joshi2024powercurve}. \tbd{The closed form and the roll chosen for the lowest wind are ours.} This is a lower bound for a kite that never climbs or turns, at the middle of the window in a uniform wind.

A pattern also has to climb. In a steady climb straight up the window with air speed $xV_a\cos\beta$ across the tether, the most force along the path is at $x=(G+\sqrt{G^2-8})/4$, or $x=0$ where that gives less:
\begin{equation}
V_\mathrm{climb} = \sqrt{\frac{2mg}{\rho_a A\cos\beta\;\max_\alpha\left[\sqrt{1+x^2}\,(C_L-C_D x)\right]}}.
\label{eq:min_climb}
\end{equation}
The balance is that of a rigid wing winched up in still air \citep{fagiano2017takeoff}, with the reel-in speed there in place of the wind along the tether here; the launch reel-in speed it reports agrees with this $x$, \tbd{whose closed form is ours}. Off the wind plane, at an angle $\theta$ from downwind, the factor $\cos\beta$ in both is $\cos^2\theta/\cos\beta$. Our quasi-steady patterns balance at each point, so $V_\mathrm{climb}$ at the lowest elevation is the wind below which the model gives no flight.

A launch starts at rest. The kite lifts off from
\begin{equation}
V_\mathrm{lift} = 1.3\sqrt{\frac{2m_0 g}{\rho_a A C_{L,\max}}},
\label{eq:lift_off}
\end{equation}
where $m_0$ is the mass of canopy, bridle and control unit. Without the factor 1.3, our estimate of the margin for gusts, control and the climb, somewhat above what an aircraft keeps over its stall \citep{faa1996}, this is the static take-off limit of a kite \citep{joshi2024powercurve,bigi2018}. SkySails adds that rough sea and bad weather might require \SI{10}{m/s}, the high bound. We keep the calm-sea figure as the base because rough sea comes with stronger wind, which clears the limit by itself. Rain, which adds weight, is not modelled.

No pattern flies below \SI{5}{m/s} at any of the six kite sizes from 400 to \SI{16000}{m^2}. A kite already up can be held still in a little less wind, from where its lift at rest carries its weight and the tether's pull, but the model keeps it up only from the first wind of its force table at which it either flies a pattern or holds still. For the \SI{1600}{m^2} kite $V_\mathrm{climb}$ is \SI{4.2}{m/s}, the launch wind \SI{5}{m/s}, a still kite first balances at \SI{4.5}{m/s} and the first pattern flies at \SI{6}{m/s}; with the table's rows \SI{1}{m/s} apart its first still kite is at \SI{5}{m/s}, so it stays up from \SI{5}{m/s} and lands between 4.5 and \SI{5}{m/s}, where it could hold still. All these limits are on the apparent wind and so count the wind of the ship's own speed; a following wind near the ship's speed cancels it. With the wind within about \SI{15}{\degree} of the bow no state of the kite pulls the ship forward; a kite kept up there is held still, and its pull backwards is charged to the ship as added resistance.

\subsection{Canopy life}
\label{sec:det_canopy_life}

A number without a citation is our estimate. \Cref{app:canopy_life_inputs} gives the data behind each input and \cref{app:canopy_life_range} the range of the estimate.

\paragraph{Decay} The strength falls as $e^{-D/D_0}$ with the UV dose $D$ and at a wear rate $r$, a strength-equivalent cloth weight lost per flight hour. The canopy ends when the strength reaches five times the skin tension $T=\SI{250}{N/m}$ at the force limit. With the margin $M=\sigma_w w/(5T)$ of new cloth, at a strength of $\sigma_w=\SI{170}{N/m}$ per g/m², within the 110 to 210 of canopy and sail cloth data \citep{porcher2006skytex_tds,ambroziak2018}, the rates add:
\begin{equation}
\frac{1}{t}=\frac{1}{t_\mathrm{UV}}+\frac{1}{t_\mathrm{wear}},\qquad
t_\mathrm{UV} = \frac{D_\mathrm{ref}}{\dot D}\left(\frac{w}{w_\mathrm{ref}}\right)^{a}\ln M,\qquad
t_\mathrm{wear} = \frac{w_\mathrm{ref}}{r}\left(\frac{w}{w_\mathrm{ref}}\right)^{b}\ln M,
\label{eq:uv_life}
\end{equation}
with $w_\mathrm{ref}=\SI{100}{g/m^2}$.

\paragraph{UV} Tests of polyester without a UV blocker, outdoors and under lamps, give $a=0.4$ and $D_\mathrm{ref}=\SI{70}{MJ/m^2}$. The dose rate $\dot D$ is \SI{0.02}{MJ/m^2} per flight hour, about 55\,\% of the yearly mean in Florida, because the kite flies day and night and is stowed out of the sun on deck when down. No source measures this dose on a kite.

\paragraph{Wear} We take $r$ from what makers state of life in service (\cref{tab:canopy_life}), after removing the UV each statement includes. The statements imply a median of \SI{0.10}{g/m^2/h}, and counting a sail's shape life in place of its structural life gives 0.19; we choose \SI{0.13}{g/m^2/h} between the two. Flutter and folds could set this value but have not been measured on a kite. The exponent $b$, how much longer a heavier cloth lasts at the same load, is our choice of 1, and its bounds change whether a lighter or a heavier cloth comes out cheaper, so we do not optimise the cloth weight.

\begin{table*}[tbp]
\centering
\footnotesize
\caption{Life statements behind the wear rate: the implied $r=w\ln M/L-w\dot D/D_0(w)$ at $a=0.4$ and $D_\mathrm{ref}=\SI{70}{MJ/m^2}$. Grades: claim, a maker's statement with no test shown; meas, a maker's measurement. An asterisk marks our estimate; $T/w$ and $\sigma_w$ are in N/m per g/m². Under $a=1$ the median is 0.15, under $a=0.1$ it is 0.09.}
\label{tab:canopy_life}
\adjustbox{max width=\textwidth}{%
\setlength{\tabcolsep}{3pt}
\begin{tabular}{@{}lllrrlrrrr@{}}
\toprule
Source & Grade & Cloth & $w$ & Life $L$ & Ends on & $T/w$ & $\sigma_w$ & $\dot D$ & $r$ \\
 & & & [g/m²] & [h] & & & & [MJ/m²/h] & [g/m²/h] \\
\midrule
\citet{flynn2015} & claim & Woven polyester sail & 250* & 3500--4000 & structure & 2.0* & 110--207 & 0.04--0.07* & 0--0.12 \\
\citet{flynn2015} & claim & same & 400* & 3500--4000 & structure & 2.0* & 110--207 & 0.04--0.07* & 0.01--0.22 \\
\citet{kubicek2026} & claim & Polyester balloon & 45 & 400--500 & not stated & 2.8* & 222* & 0.01--0.03* & 0.22--0.30 \\
\citet{kubicek2026} & claim & Polyester balloon & 70 & 800--1000 & not stated & 2.9* & 222* & 0.01--0.03* & 0.16--0.23 \\
\citet{kubicek2026} & claim & Polyester balloon & 77 & over 1000 & not stated & 2.6* & 222* & 0.01--0.03* & 0.18--0.21 \\
\citet{supair2016} & meas & Tandem paraglider, fibre unstated & 38--45* & 800--850 & porosity & 1.6* & 200* & 0.05--0.09* & 0.07--0.14 \\
\citet{lynn2007} & claim & Maxi show kite & 50* & 1800--2400* & porosity & 0.6* & 200* & 0.05--0.09* & 0--0.07 \\
\bottomrule
\end{tabular}
}
\end{table*}

\paragraph{What no source settles} The model ends a canopy on lost strength, while flex, folds, flutter and abrasion leak the cloth and UV takes strength with little leak \citep{maslanka2023,maslanka2023thesis}; no leak limit is published for a kite. Retiring a kite at the 27 to 70\,\% of new strength used for balloon cloth and aircraft fabric gives the low end of the range. Launches are not counted, although a canopy sees about \num{2500} launch and landing cycles at the 1.5 hours a flight SkySails reported \citep{skysails2018}. The decay does not know salt, heat or a ship's handling.

\subsection{Cloth tension at the force limit}
\label{sec:cloth_check}

This check feeds no cost (\cref{fig:canopy_tension}). The tether force $F$ is carried by a pressure $p=F/A$ over the canopy, and a skin bulging to radius $R$ between ribs carries $T=pR$ per unit width. The force limit caps the peak at $p=\SI{500}{N/m^2}$. We state one radius, $R=\SI{0.5}{m}$; the cell equilibrium of \citet{wang2018} gives \SI{0.41}{m} on the SkySails kite (\cref{app:cloth_tension}), and cells of that width on the chord of the \SI{3200}{m^2} kite bulge to \SI{1.2}{m}.

The cloth carries \SI{250}{N/m} at the peak. With the pressure $p(1+1/C_L)$ on the top skin near the leading edge \citep{wang2018}, the margin of new cloth of \SI{200}{g/m^2} falls from 136 to about 60 times, and at a radius of \SI{1.2}{m} to 55 times. The model leaves out gusts, ship motion, launch loads and flutter, and no test of canopy cloth covers the range of \SI{71}{N/m} at 44 cycles an hour.

\begin{figure*}[tbp]
\centering
\plot[\textwidth]{canopy_tension}
\caption{From tether force to cloth tension on the \SI{3200}{m^2} kite at $R=\SI{0.5}{m}$. (a) The force balance at the force limit. (b) Tether force and area-averaged cloth tension over one optimal cycle in \SI{12}{m/s} of apparent wind, \SI{90}{\degree} off downwind, with the strength of new cloth of \SI{200}{g/m^2}.}
\label{fig:canopy_tension}
\end{figure*}

\subsection{VTOL kite}
\label{sec:det_rigid}

\paragraph{Why it hovers} A rigid wing's static lift alone asks for 17 to \SI{25}{m/s} of relative wind at the wing loadings here, and other launches need far more room than a deck has or have been shown at small scale only (\cref{app:vtol_data}). Part~II of the report on the one offshore hover calls launch and landing robust only onshore \citep{echeverri2020part2}. SkySails authors expect heavy rigid structures to be hard to stow on a ship \citep{paulig2013}; the model charges a perch, its deck space and the hover but nothing for a harder stow.

\paragraph{Mass and scale} The mass of a wing with its propellers goes with span to a power $\kappa$ at a fixed aspect ratio,
\begin{equation}
m = \SI{1852}{kg}\left(\frac{A}{\SI{54}{m^2}}\right)^{\kappa/2},
\label{eq:rigid_mass}
\end{equation}
with $\kappa=2.75$ and bounds of 2.2 and 3.3 (\cref{fig:rigid_mass}). The anchor is MX2, and the exponent lies between the two values Wing~7 of Makani \citep{vanderlind2013} gives to it. This is not a fit, since all three points come from one company and MX2 was never built (\cref{app:vtol_data}). The M600 came out 40 to 70\,\% over its targets \citep{echeverri2020part1}, so the risk is upward, and the high bound of the reference mass adds 40\,\%. The two wings actually built go about as span cubed, so the size study also runs $\kappa=3$. The tether is sized as on MX2, and its drag coefficient stays at 0.029 at every size.

\paragraph{Polar and propeller drag} The lift rises on the slope of a finite wing of aspect ratio 12.5 to a highest lift coefficient of 2.0, and the drag is a profile drag of 0.03 plus the induced drag at a span efficiency of 0.9, which gives the MX2 drag coefficient of 0.123 at its lift coefficient of 1.81 \citep{echeverri2020part1,makani2020code}. An idle propeller windmills, at the drag coefficient of 0.022 of that release's propeller table, from simulation and not from flight (\cref{app:vtol_data}). The largest open risk is the parasitic drag: a flight glide of the M600 gave an offset of 0.075 against the 0.005 assumed for MX2, carried in the high bound of the profile drag. Hovering needs nearly the same drives, so we claim no saving on the drive mass.

\paragraph{Loops} Almost no steered sweep or figure eight is flyable for a wing this heavy. It flies circles on a cone about an axis that the search places for each wind and course, on a radius of at least 3.46 spans, the \SI{90}{m} of MX2. On the larger loops and longer tethers the wings that flew needed, the \SI{54}{m^2} kite flies from \SI{9}{m/s} and not \SI{7}{m/s} (\cref{app:vtol_data}). The balance of \cref{eq:kite_drag,eq:kite_lift} is solved at 36 points of the loop, each counting for the time it takes, with no inertia between points; the sign of that error at this scale is unknown. A loop is flown only if the mean airspeed pays for the climb to its top above the stall speed,
\begin{equation}
\overline{v_a}^2 - g\,\Delta h \geq \frac{2mg}{\rho_a A C_{L,\max}},
\label{eq:loop_top}
\end{equation}
and $\Delta h$ grows with the span while the stall speed grows with $m/A$, so this limit sets the lowest wind of a large wing whatever its mass law, as \citet{vanderlind2013} argues for Wing~7. A margin above the stall speed would raise the lowest wind of the \SI{54}{m^2} kite by \SI{1}{m/s}; we keep the stall speed because our model already lies above Wing~7's published curves. The lowest wind is the lowest in the table at which the kite flies on any course, and at least $V_\mathrm{cross}$ of \cref{eq:min_cross}. It flies up to \SI{25}{m/s} of apparent wind, the limit of the soft kite, for lack of a published one.

\paragraph{Several kites} Several kites on one ship share one loop, spaced evenly around one cone, so that their tethers cannot cross. They act like the blades of a wind turbine on the disc of that loop, and momentum theory gives each of $n$ kites the share
\begin{equation}
\left(\frac{4+C}{4+nC}\right)^2
\label{eq:loop_share}
\end{equation}
of the force of a lone kite, with $C$ the force of a lone kite over the dynamic pressure and the area of the disc. We use this rather than the array factor of \cref{sec:det_deck}, and no flight of several tethered kites on one loop supports it. Each kite brings its own tether, winch, perch and propeller drive, \SI{20}{t} of station and \SI{8}{m} of deck that blocks cargo there, and perches stand a span apart, which limits the fleet on a short ship.

\paragraph{Launch, landing and sea} A kite hovers for \SI{6}{min} over one launch and one landing, with no force on the ship, on electricity from the ship at \SI{0.20}{USD/kWh}, at \SI{873}{kW} for the \SI{54}{m^2} kite from momentum theory (\cref{app:vtol_data}). It needs no wind to lift off and no turn, and we charge it no retry. No source gives limits on wind or sea for its launch and landing, so we estimate an apparent wind of \SI{17}{m/s}, with bounds of 13 and \SI{20}{m/s}, and a significant wave height of \SI{3}{m}, with bounds of 1.5 and \SI{4}{m}. Aloft, like the soft kite, it has no limit on the sea.

\paragraph{Prices and life} No VTOL kite has a public price. We build one from the cost model of \citet{joshi2025}, without margin, so the market price factor of the paper applies. The drive, tether, perch and upkeep follow from \citet{bauer2018}, \citet{echeverri2020part1} and our estimates (\cref{app:vtol_data}). We take 5 in \num{10000} as today's estimate of the loss per launch and landing, and USD~2.3 million results for the \SI{54}{m^2} kite with its tether, winch and perch.

\subsection{Routing search}
\label{sec:det_routing}

Every route uses the same corridor both ways. When a unit is out of service (\cref{sec:cost_model}), those still working of several VTOL kites fly on, and where a kite flies over a row of rotors the kite and each rotor fail on their own.

The dynamic programme of \cref{sec:routing} keeps the cheapest arrival per state, carrying its exact arrival time rather than the time of its bin, and traces back the best final state, for a kite the best with the kite down. A ship at \SI{15}{kn} sails a quarter of a stage or less in one bin, the condition of \citet{wendling2025} that keeps the ship from skipping a stage within one time step.

Where the engine cannot give the thrust for the step's speed, the ship slows there to the highest speed it can hold. Which speeds it can hold depends on the kite's state when the step begins: with the kite down, they are those it holds with the kite stowed, and a flying kite sets the speed only where none of those holds; with the kite up, they include speeds that only a flying kite makes. A ship that launches during a step thus keeps the speeds of its stowed kite until the next step, and a ship whose kite cannot launch can still sail any step it holds a speed on with the kite stowed. Where no speed down to \SI{4}{kn} holds, the step is removed.

A coarser backward programme first estimates the cost from each node and aloft state to the destination, on time knots \SI{24}{h} apart. The forward programme skips any state or edge whose cost so far plus this estimate exceeds the estimate at departure by more than 3\,\%, or 20\,\% for a ship that can fly a kite. If its best passage ends above that limit, or it finds none, the search is repeated at 6\,\%, or 40\,\% with a kite, and then over all states. The estimate interpolates linearly in time between its knots, so it cannot time a passage to a spell of good wind: on the worst of 36 legs it reads 23\,\% high for a kite and 22\,\% for rotors with knots \SI{48}{h} apart, and 15\,\% and 10\,\% at \SI{24}{h}, while finer weather samples leave that unchanged and more speeds change it by about two points. Averaged over those legs it reads high by 6.5\,\% for a kite, 5.6\,\% for rotors, 2.3\,\% for wing sails and 1.1\,\% without a rig at \SI{48}{h}. How far the cost to go departs on average from a straight line between knots ranks the rigs nearly the same way, at 1.9\,\% for a kite, 2.4\,\% for rotors, 1.3\,\% for wing sails and 0.7\,\% without a rig, but it explains less than half of the difference between legs; the rest, which leaves a kite reading 2.6\,\% high on average even with knots \SI{6}{h} apart, we have not traced. Knots \SI{12}{h} and \SI{6}{h} apart read closer still but lost one rotor case 0.44 points of cost saving, so the spacing stops at \SI{24}{h}. Against a search of all states on 24 cases, two ships on three routes, bare and with each rig, at stages \SI{360}{nmi} and \SI{180}{nmi} apart, these settings leave the kite's and the wing sails' cost saving no lower, the rotors' at most 0.01 points lower and the ship without a rig unchanged, at about 30\,\% more run time than knots \SI{48}{h} apart with kite margins of 10\,\%, then 20\,\%, which lost up to 0.27 points; because a final pass refines whichever path the search found, a rig's saving can also come out up to 0.08 points higher than with all states searched. Kite margins of 10 or 15\,\%, doubled on the retry, lose up to 0.04 points, and of 7\,\% up to 0.13. Taking the search limit from the best cost on the next coarser corridor instead lost up to 0.94 points on the kite at up to twice the run time. The estimate can lie on either side of the true cost, so the result is not exact in principle. Two cheaper ways of discarding states, dropping an arrival later and dearer than another and stopping at the first speed that costs more, lost 2 to 5 points of cost saving, because a later, dearer arrival often meets better weather.

The search tries the speeds and time bins of \cref{tab:numerics} and lets an edge move at most two lateral nodes. A final pass repeats the programme along the path found, at all 49 speeds of the table and in time bins of \SI{0.25}{h}, and its passage is the result. The speeds are a share of each ship's own range, so that a larger, faster ship is not routed on more speeds. Weather stays sampled at least every \SI{3}{h} along an edge, because \SI{6}{h} raises the kite's cost saving by up to 0.6 points. \Cref{app:routing_convergence} gives the convergence of these settings.

The worth of searching with the rig (\cref{sec:routing}) is measured on the 12 departures of the \SI{140}{m} e-methanol ship and the Kamsarmax on VLSFO on the routes to Guadeloupe, New York and Tubarão, with 1600 and \SI{3200}{m^2} of kite, 1500 and \SI{3000}{m^2} of wing sails and 800 and \SI{1600}{m^2} of rotors, outbound and home. The cost is the one the router minimises, without the rig's capital, and every route of a case uses the time price of the ship with the rig. The path found without a rig, sailed in a final pass at all speeds, comes out cheaper than the search's own on 12 of the kite's 144 passages, on 3 of the wing sails', by at most 0.2\,\% of the cost of the ship without a rig, and on 1 of the rotors', by 1.0\,\%.

\begin{table*}[tbp]
\centering
\footnotesize
\caption{This router and the routers it is compared with, each as its publication describes it. Speed is the ship's speed through the water or, where marked, the engine power that sets it.}
\label{tab:routers}
\begin{tabular}{@{}p{2.4cm}p{3.1cm}p{2.2cm}p{2.7cm}p{3.8cm}@{}}
\toprule
Router & Search & Speed & Objective and arrival & Rig, and how the search treats it \\
\midrule
This paper & forward dynamic programming on a corridor with time bins & chosen per step & money: fuel, kite wear and launches, time at a price; arrival free & kite, wing sails, rotors; kite up or down in the search state, launch and landing charged \\
\citet{wei2012} & forward dynamic programming on stages, lateral points and time & chosen per step & fuel at a given arrival time & none \\
\citet{mason2023} & dynamic programming over position and speed, arrival time rounded & chosen per stage, 8 to \SI{14}{kn} & fuel, arrival limit & rotors, set by a rule in the performance model \\
\citet{wendling2025} & dynamic programming over position and time & engine power chosen per arc & front of time against energy & wing sails, in speed tables made before the search \\
\citet{dhainaut2025} & gradient-based optimal control, to a local optimum & continuous & fuel, maximum travel time & rotors, spin and position trimmed continuously, no on or off \\
\citet{wang2021} & genetic algorithm seeded by shortest paths and dynamic programming & eight engine power levels per leg & fuel and closeness to an arrival time & none \\
\citet{mannarini2024} (VISIR-2) & shortest path in time-varying weather, one cost per point & engine load set per run & distance, time or CO$_2$; arrival free & sailing boat, rig fixed before the search \\
\citet{marzi2024} & isofuel, which extends each route by an equal amount of fuel & constant for the voyage & fuel & none \\
\citet{bentin2016,bentin2018} & A$^*$, a shortest-path search guided by an estimate & \SI{13}{kn} held & energy, maximum travel time & rotors (2016); kite, rotors, DynaRig (2018), as force models; no launch or landing \\
\citet{leskinen2024} & shortest path on a variable grid & set by the wind, electric motor below that & least time & kite, in a table of ship speed; no launch or landing \\
\bottomrule
\end{tabular}
\end{table*}

\subsection{Forecast router}
\label{app:forecast_router}

The forecast router re-plans at the first node of the route 24~hours or more after the last plan, on forecasts that were actually issued: the control run of NOAA's GEFSv12 reforecast \citep{guan2022} for the wind at 10~m, and the WAVEWATCH~III wave reforecast driven by it \citep{campos2024} for the height, period and direction of the waves. Both start once a day, reach 16~days ahead and cover 2000 to 2019. The base plan uses one run of the forecast model. The mean of its ensemble scores better but smooths away the storms beyond a week \citep{campos2024}, so whether a plan on it is the better one has to be tested. Three further plans use the control run and four perturbed members of the same reforecast: the mean of the five, each member in turn with the route that costs least on average over the five, a member on which the ship cannot hold the plan's speed counting the hours it loses slowing, and a margin for wind that may not come, the lower quartile of the members' wind speeds along their mean wind and the upper quartile of their wave heights. A further plan on the control run keeps the kite inside limits 20\,\% tighter than the real ones, its lowest flying and launch winds and the extra launch wind in a sea raised and its highest flying and launch winds lowered by that share, its landing limits unchanged; the kite then sails against its real limits. We remove the mean difference between each forecast field and ERA5 for every grid point and lead time, estimated from the forecasts issued in the other two years, as reforecasts are calibrated \citep{guan2022}; the wave direction is not corrected. The error that remains comes from the forecast model itself, so it is correlated over large areas and over days as real forecast error is, which a random error laid over ERA5 would not be without modelling the weather. The ship sails the first day of each plan in the real weather on the plan's course and at its speed, and every passage is scored in ERA5. Its kite launches and lands wherever that pays in the weather the ship meets, within the launch, landing and night limits. \tbd{Given the real weather as its forecast for 1000 days ahead, this router reproduces the cost of the base case to within 0.01\,\% for the ship without a rig, the wing sails and the rotors and 0.3\,\% for the kite, so the loss on forecasts is not an artefact of re-planning; from a draft run.}

Two plans change how the router plans rather than what it plans on. The first weighs the wind error: at each point it takes the wind speed as normally distributed about the forecast, with a standard error of the root-mean-square error of the corrected vector wind at that lead along this route over $\sqrt{2}$, the error growing from \SI{1.25}{m/s} at six hours to \SI{4.0}{m/s} at four days (\cref{fig:forecast}). It plans on the expected fuel and the chance that the kite stays up, over the forecast wind speed and $\sqrt{3}$ standard errors either side, weighted 2/3, 1/6 and 1/6 as in a three-point Gauss rule; where the kite cannot stay up in one of these winds the ship burns there the fuel it burns with the kite stowed. It needs no run of the ensemble. Its search is steered, as every plan's is, by a backward estimate of the cost of the rest of the passage, which it takes on the forecast without weighing the wind error, and its final pass bins time by the hour instead of the quarter hour (\cref{tab:numerics}). The second plans on a grid fine near the ship and coarse beyond. A ship re-plans after about a day, so it sails only the first edges of each plan, while the forecast is worst far ahead, where the plan is thrown away before it is reached. These plans sail a route grid of stages \SI{240}{nmi} and rows \SI{60}{nmi} apart, on which an edge reaches 11 headings, from straight on to \SI{51}{\degree} off the corridor, against three on the draft's grid. The ship plans on that grid for the first \SI{240}{nmi} ahead and beyond them on stages about \SI{720}{nmi} apart, on every eighth row, whose edges reach as many rows as their length allows, with the weather sampled every 12~hours; the track along the coast stays whole. \tbd{In an earlier draft run, planning on the route grid all the way kept $11 \pm 5$ points more of the kite's cost saving on one forecast run than a fine zone of \SI{480}{nmi} with every fourth row beyond it, $6 \pm 3$ more of the rotors' and $20 \pm 8$ more of the wing sails', at seven times the run time; the lighter zone used here keeps as much of each rig's saving as that one within the error on the same 12 departures, weighing the wind error, at a third of its run time. The grid's coarseness thus leaves the forecast plan keeping somewhat less than a finer one would. In the run of \cref{tab:forecast_plans} a round trip of the kite ship takes 44~s to plan on one forecast run on the draft's grid, 146~s on several members, 76~s on one run on the grid fine near the ship and 61~s on that grid weighing the wind error, and of the ship with rotors 40, 105, 71 and 47~s, on four threads of one machine shared with other work, which slows every plan, so the ratios between plans are firmer than the seconds.}

The re-planning study (\cref{sec:res_robustness}) replaces the wind reforecast with the control run of the operational forecasts of NOAA's Global Ensemble Forecast System in the version of those years, GEFSv11, one model version before the GEFSv12 of the reforecast, issued four times a day and held in NOAA's open archive (\url{https://registry.opendata.aws/noaa-gefs}), corrected for its mean difference from ERA5 in the same way, each issue reaching the ship six hours after it starts; the waves stay on the daily wave reforecast. Where an issue is missing, 12 of the 4380, each incomplete or unreadable in NOAA's archive, the ship uses the one before. Between plans a ship that adapts chooses each step's speed, or also its next node, by the objective of that step in the plan's forecast corrected by the difference it meets at its own position, of the wind vector and the wave height, held over the step, plus an estimate of the objective of the rest of the passage from the node the step ends on, in the forecast the plan used. It is the backward estimate the route search itself uses, a coarser search of the corridor that can lie on either side of the true value, because a ship adapting its next node can end a step off the plan, where the plan has no remaining cost of its own.

The forecast wind is averaged over each cell of the study's ERA5 wind grid and the forecast waves are taken at the points of its wave grid. A day's forecast reaches the ship at 06~UTC. Beyond its last lead the router sees the weather of the same date in another year, and near the coast, where the forecast holds no waves, the waves of another year. Where a wave forecast is missing, on one day of the control run and on five of the 2736 of the perturbed members, the ship uses the forecast of the day before. Along the route to Guadeloupe the error of the corrected wind and wave height grows to near that of another year's weather by ten days ahead (\cref{fig:forecast}).

At the end of each step the kite is up or down, whichever costs less on that step in the real weather plus what the rest of the plan costs from there on the forecast, so the choice counts the launches it costs later. A kite held to its plan instead keeps the state the plan gave each point, launching or landing only where the real limits allow, which shows what flying the weather met is worth. Where the plan's speed can be held only with a kite the real seas do not let it launch, the ship slows to the highest speed it holds with the kite as it is. Where the real weather leaves no route from the ship's position, or stops a step of its plan at every speed, the ship holds there on its hotel load with the kite stowed and re-plans on the next forecast, 6~hours later. After 72~hours at one point its passage fails. \tbd{On the 24 passages of the 12 departures the kite's state differs between the two rules on 176 of the 4475 samples. Step by step, each step taken from the state of the kite that chooses, following the plan costs 2.3\,\% more in fuel and rig on the 21 passages where it can sail every step at that kite's speeds, and on the other three the plan's kite state cannot hold those speeds on some step. The cost of a kite held to its plan through whole passages is in \cref{tab:forecast_plans}.}

\subsection{How the router picks speed}
\label{app:speed_choice}

A rig that gave the same thrust at every speed would not change the best speed: in the cost per mile $(R(V)-F_x)/\eta + c/V$, for a resistance $R$, an efficiency $\eta$ and a cost $c$ per hour of time, a constant $F_x$ drops out of the derivative. \Cref{tab:speed_choice} shows what moves it instead, on the \SI{140}{m} ship in steady wind (\cref{sec:res_speed}). The kite moves it by the speed its pull alone can reach, since up to there the ship burns only its hotel load and sailing faster costs nothing. The rotors' best speed rises mainly because their force grows with speed. The wing sails move it by a knot or less, mostly by lightening the load on the propeller. The capital and wear of the rig shift it by about \tbd{0.1}~kn.

\begin{table}[tbp]
\centering
\caption{Cost-optimal speed [kn] of the \SI{140}{m} ship in steady wind, median over winds of 6, 10 and \SI{14}{m/s} at 60, 90, 120 and \SI{150}{\degree}, with the reference rigs: with the rig's capital and wear left out of the cost, with its force held at its value at the speed of the ship without a rig, and the speed at which that held force alone moves the ship. \tbd{Draft run.}}
\label{tab:speed_choice}
\adjustbox{max width=\columnwidth}{\InputIfFileExists{sections/speed_choice_table}{}{\tbd{Speed table pending: run the speed\_choice\_table step.}}}
\end{table}
